\documentclass[12pt,a4paper]{article}

% =========================================
% PAGE SETUP
% =========================================
\usepackage[
    a4paper,
    margin=25.4mm
]{geometry}


% =========================================
% BASIC PACKAGES
% =========================================
\usepackage{graphicx}       % Images
\usepackage{amsmath}        % Mathematics
\usepackage{amssymb}        % Mathematical symbols
\usepackage{booktabs}       % Professional tables
\usepackage{float}          % [H] figure/table positioning
\usepackage{caption}        % Caption customization
\usepackage{hyperref}       % Hyperlinks and references
\usepackage{ragged2e}       % Text alignment
\usepackage{setspace}       % Line spacing
\usepackage{hyperref}       % Clickable links and references
\usepackage{array}          % Array for tables (column types for timeline/milestone tables)
\usepackage{algorithm}      % Float environment for pseudocode/algorithm listing

% Column types for proposal tables (wrapping, left/center aligned)
\newcolumntype{L}[1]{>{\RaggedRight\arraybackslash}p{#1}}
\newcolumntype{C}[1]{>{\Centering\arraybackslash}p{#1}}

% =========================================
% APA REFERENCING
% =========================================

\usepackage[
    style=apa,
    backend=biber
]{biblatex}

\addbibresource{references.bib}


% =========================================
% HYPERLINK SETTINGS
% =========================================
\hypersetup{
    colorlinks=true,
    linkcolor=black,
    urlcolor=black,
    citecolor=black,
    pdftitle={Research Proposal}
}

% Table and figure caption formatting
\DeclareCaptionFont{captionnine}{\fontsize{9pt}{11pt}\selectfont}

\captionsetup[table]{
    font=captionnine,
    justification=raggedright,
    singlelinecheck=false,
    labelsep=colon
}

\captionsetup[figure]{
    font=captionnine,
    justification=centering,
    singlelinecheck=false,
    labelsep=colon
}


% =========================================
% DOCUMENT
% =========================================
\begin{document}

\begin{center}

{\fontsize{18}{22}\selectfont\textbf{RESEARCH PROPOSAL}}

\vspace{0.5cm}

{\fontsize{18}{22}\selectfont\textbf{Adaptive Gamification with Real-Time Harm Mitigation for Personalised Learning: A Quasi-Experimental Proposal}}

\vspace{1cm}

{\fontsize{12}{14}\selectfont\textbf{Group Members}}

\vspace{0.3cm}

{\fontsize{12}{14}\selectfont
1. IMRAN HARIS BIN MOHD AZLI (33\%) -- 252UC254P6\\
2. LEE CHONG CHUN (34\%) -- 252UC254YW\\
3. SHAM CHALAN (0\%) -- Unknown\\ % This student did not answer to emails and replies sent to them
4. YAW YI WEOI (33\%) -- 251UC250TR
}

\vspace{0.8cm}

{\fontsize{12}{14}\selectfont
\textbf{Programme:} Bachelor of Computer Science\\
\textbf{Course:} CPT6123-RESEARCH METHOD IN COMPUTER SC\\
\textbf{Lecturer:} Jamal Arshad\\
\textbf{Semester:} T2620\\
\textbf{Submission Date:} 27\textsuperscript{th} Sep 2026
}

\end{center}


% =========================================
% ABSTRACT
% =========================================

\vspace{0.5cm}

\noindent
\textbf{Abstract}

\vspace{0.2cm}

\begin{justify}

Educational gamification now produces well-documented engagement and achievement gains alongside equally well-documented harms, including competitive leaderboard anxiety, undermining of intrinsic motivation, and superficial point-chasing, yet most platforms still deploy one fixed configuration to every learner with no mechanism to detect or constrain harm as it occurs. This proposal specifies an adaptive, harm-aware gamification controller that estimates a learner state of engagement and anxiety from platform interaction traces, computes a quantitative harm score against a defined threshold, and selects, retains, or downgrades gamification elements under that constraint, together with a semester-long quasi-experimental evaluation comparing the controller with static gamification and non-gamified baselines. The study is expected to yield a validated real-time learner-state model, an operational harm metric, and a set of element-configuration rules that preserve the moderate-to-large engagement and achievement effects reported in the literature while reducing anxiety and motivational degradation, thereby providing educators with an evidence-based and ethically accountable basis for deploying gamification at scale.

\vspace{0.3cm}

\textbf{Keywords:} Adaptive Gamification; Harm Mitigation; Learner Modelling; Quasi-Experimental Design; Student Engagement

\end{justify}




% =========================================
% TABLE OF CONTENTS
% =========================================

\newpage

\tableofcontents

\newpage



% =========================================
% MAIN CONTENT
% =========================================

\begin{justify}

% =========================================
% 1. INTRODUCTION
% =========================================

\section{Introduction}

% Establish the broader CS field (adaptive/intelligent learning systems) and explain
% why adaptive gamification with harm mitigation matters in the current context.
% Include: background, relevance, brief roadmap of the proposal.

[Introduction: broader context of intelligent adaptive learning systems and why
real-time, harm-aware gamification matters now.]


% =========================================
% 2. LITERATURE REVIEW
% =========================================

\section{Literature Review}

This proposal builds on the group's preceding critical research review of gamification in education, which synthesised recent peer-reviewed studies published between 2022 and 2025 spanning theoretical analysis of motivation, meta-analytic evidence synthesis, systematic documentation of negative effects, and classroom implementation in higher education, artificial intelligence education, web development, and vocational training. Rather than repeating that review article by article, the literature is organised below thematically around the three research questions of this proposal, so that the theoretical foundation, the documented harms, and the available evaluation evidence can be read directly against the design problem this proposal addresses.

\subsection{Adaptive and Personalised Gamification (Link to RQ1)}

[Thematic synthesis of adaptive/learner-model gamification studies.]

\subsection{Documented Negative Effects and Harm-Mitigating Design (Link to RQ2)}

[Thematic synthesis of negative-effect and anti-addiction/harm-mitigation studies.]

\subsection{Evaluation Designs and Evidence for Gamification Effectiveness (Link to RQ3)}

[Thematic synthesis of experimental/quasi-experimental and meta-analytic evidence.]

\subsection{Gap Identification}

Three gaps follow from the preceding themes, and they map directly onto the research questions of this proposal. In relation to RQ1, the literature affirms that learners respond differently to identical elements and repeatedly recommends personalisation, yet it provides no real-time learner-state model and no validated feedback loop between an estimated state and element selection. In relation to RQ2, harms are well documented and attributable to particular design choices, yet no operational harm metric exists that a system could compute online and enforce as a constraint on deployment. In relation to RQ3, evidence of effectiveness is substantial but derives from short, non-randomised, single-condition comparisons that measure benefits without commensurate measurement of harms, leaving the central practical question of whether an adaptive, harm-aware configuration preserves attainment gains while reducing documented negative effects unanswered. The proposal addresses these three gaps jointly, because a harm metric is actionable only if learner state can be estimated in real time, and a state model is meaningful only if its output constrains element selection and is subsequently evaluated against both benefit and harm criteria. The problem statement below formalises this transition from context to gap to proposed solution.


% =========================================
% 3. PROBLEM STATEMENT
% =========================================

\section{Problem Statement}

% TODO: CONTEXT - GAP - SOLUTION
% Use the Context--Gap--Solution method. Justify with data/gaps from the Assignment 1
% literature review.

\subsection{Context}

[What is currently known about adaptive gamification and its negative effects.]

\subsection{Gap}

[The specific technical gap: no closed-loop, harm-aware adaptation with a quantitative harm metric.]

\subsection{Solution}

This proposal specifies an adaptive, harm-aware gamification controller that addresses the three deficiencies as a single closed loop. A learner-state model fuses engagement and anxiety signals from the platform interaction stream to produce a real-time state estimate; a harm function maps that estimate and the currently deployed elements to a scalar harm score, which is compared against a threshold $\tau$; and a selection policy proposes the next element configuration, downgrading competitive and reward-heavy elements whenever the harm score exceeds $\tau$ and otherwise selecting the configuration predicted to sustain engagement. Element selection is therefore constrained rather than merely optimised, so documented negative effects become an explicit design variable instead of an unmonitored side effect. The proposed system is to be evaluated through a quasi-experimental design over one full semester, comparing the adaptive harm-aware condition with a static gamification condition and a non-gamified baseline using balanced benefit-and-harm metrics, thereby producing the first evidence on whether real-time harm mitigation preserves the engagement and achievement gains reported in the literature while reducing anxiety and motivational undermining.


% =========================================
% 4. RESEARCH QUESTIONS
% =========================================

\section{Research Questions}

\vspace{0.2cm}

\begin{enumerate}
    \item How can learner state (engagement and anxiety) be modelled in real time to drive adaptive selection of gamification elements in an online learning platform?
    \item Which gamification design configurations minimise documented negative effects -- anxiety, motivational undermining, and superficial learning -- while preserving the engagement and achievement gains reported in the literature?
    \item Does an adaptive, harm-aware gamification system improve learning outcomes and motivational quality compared with static gamification and non-gamified baselines over a full semester?
\end{enumerate}


% =========================================
% 5. RESEARCH OBJECTIVES
% =========================================

\section{Research Objectives}

\begin{enumerate}
    \item \textbf{RO1 (links to RQ1):} [SMART objective -- e.g., design and specify a real-time learner-state
    model that fuses engagement and anxiety signals and selects gamification elements adaptively.]
    \item \textbf{RO2 (links to RQ2):} [SMART objective -- e.g., define a quantitative harm metric and a set of
    element-configuration rules that constrain deployment to minimise documented negative effects.]
    \item \textbf{RO3 (links to RQ3):} [SMART objective -- e.g., evaluate the adaptive harm-aware system against
    static gamification and non-gamified baselines over one semester using defined benefit-and-harm metrics.]
\end{enumerate}


% =========================================
% 6. RESEARCH METHODOLOGY
% =========================================

\section{Research Methodology}

\subsection{Research Design and Approach}

% Choose: Experimental, Quasi-experimental, Empirical, or Comparative -- with justification.
% Include a pipeline overview (conceptual flow of data -> processing -> output).

[Justified quasi-experimental design; high-level pipeline described in text.]

\subsection{Technical Approach}

% Step-by-step breakdown of the proposed algorithm/model/framework.
% Model selection justification, algorithmic innovation (novel vs. repurposed),
% optimisation strategy. MUST include pseudocode or a high-level architecture
% description compatible with LaTeX.

\begin{algorithm}[H]
\caption{Adaptive harm-aware gamification controller (pseudocode outline)}
\label{alg:controller}
\begin{flushleft}
\fontsize{10pt}{12pt}\selectfont
\textbf{Input:} learner interaction stream $\mathbf{x}_t$, harm threshold $\tau$ \\
\textbf{Output:} gamification element configuration $e_t$
\end{flushleft}
\hrulefill
\vspace{0.2cm}
\begin{enumerate}
    \item Estimate learner state $\hat{s}_t = f_\theta(\mathbf{x}_t)$   % engagement + anxiety model
    \item Compute harm score $h(\hat{s}_t, e_{t-1})$
    \item \textbf{if} $h > \tau$ \textbf{then} $e_t \leftarrow$ downgrade competitive elements
    \item \textbf{else} $e_t \leftarrow \pi_\theta(\hat{s}_t)$   % adaptive element selection policy
    \item Observe outcome; update $\theta$   % optimisation / learning step
\end{enumerate}
\hrulefill
\end{algorithm}

\subsection{Data and Tools}

% Datasets (provenance: public benchmark vs. platform-collected), key tools
% (e.g., Python, scikit-learn/PyTorch, Moodle/LMS logs), justification of size/noise/features,
% ethical considerations (human subjects, consent, privacy, bias mitigation).

[Dataset sources, platform instrumentation, software tools, and ethical considerations.]

\subsection{Evaluation Criteria}

% Clear measurable metrics (learning gain, engagement rate, motivational-quality score,
% harm metric values), experimental setup (hardware/software), baselines and ablation
% (static gamification, non-gamified control; element-level ablation).

[Metrics, baselines, ablation plan, and experimental setup.]


% =========================================
% 7. EXPECTED OUTCOMES
% =========================================

\section{Expected Outcomes}

% Anticipated findings, significance to the field, deliverables.
% Must be specific and measurable.

[Anticipated findings, contributions, and deliverables.]


% =========================================
% 8. PROJECT TIMELINE
% =========================================

\section{Project Timeline}

% Phased structure with critical milestones and checkpoints (Gantt-style table).

\begin{table}[H]
\caption{Phased project timeline with milestones}
\label{tab:timeline}
\fontsize{10pt}{12pt}\selectfont
\begin{tabular}{|L{3.2cm}|L{4.5cm}|L{3cm}|L{3cm}|}
\hline
\centering\textbf{Phase} & \centering\textbf{Activities} & \centering\textbf{Duration} & \centering\textbf{Milestone}
\tabularnewline
\hline
1. Requirement \& design & Gap analysis, system specification & Days 1--2 & Design document
\tabularnewline
\hline
2. Modelling \& prototyping & Learner-state model, controller prototype & Days 3--6 & Working prototype
\tabularnewline
\hline
3. Pilot \& refinement & Pilot run, metric calibration & Days 7--8 & Pilot report
\tabularnewline
\hline
4. Full evaluation & Baselines, ablation, analysis & Days 9--12 & Results \& write-up
\tabularnewline
\hline
\end{tabular}
\end{table}


% =========================================
% 9. CONCLUSION
% =========================================

\section{Conclusion}

\subsection{Summary of Contributions}

[Recap the proposed solution and its expected contribution -- no new information.]

\subsection{Final Remarks}

[The take-home message: harm-aware adaptive gamification as a responsible next step
for technology-enhanced learning.]


\end{justify}

% =========================================
% REFERENCES
% =========================================

\newpage

\phantomsection
\addcontentsline{toc}{section}{References}

\printbibliography[
    title={References}
]





% =========================================
% APPENDIX A - INDIVIDUAL ARTICLE REVIEW
% =========================================

\newpage

\phantomsection
\addcontentsline{toc}{section}{Appendix A -- Individual Article Review}

\begin{flushleft}
    {\fontsize{14pt}{17pt}\selectfont
    \textbf{APPENDIX A -- INDIVIDUAL ARTICLE REVIEW}}
\end{flushleft}

\vspace{0.5cm}

Each student should complete the following table for their
\textbf{minimum 3 research articles} before the group combines the findings.

\vspace{0.5cm}

\textbf{Student:} [Name / Student ID]


\begin{table}[H]

\caption{Individual article review}
\label{tab:individual-article-review}

\fontsize{10pt}{12pt}\selectfont

\noindent\resizebox{\textwidth}{!}{%
\begin{tabular}{|L{2.5cm}|L{2.5cm}|L{2cm}|L{2.5cm}|L{2.5cm}|L{2.5cm}|L{2.5cm}|L{2.5cm}|}
\hline

\centering\textbf{Article} &
\centering\textbf{Research Problem} &
\centering\textbf{Objective} &
\centering\textbf{Methodology} &
\centering\textbf{Dataset/Sample} &
\centering\textbf{Key Findings} &
\centering\textbf{Limitations} &
\centering\textbf{Research Gap}
\tabularnewline
\hline

&  &  &  &  &  &  & 
\tabularnewline
\hline

&  &  &  &  &  &  & 
\tabularnewline
\hline

&  &  &  &  &  &  & 
\tabularnewline
\hline

&  &  &  &  &  &  & 
\tabularnewline
\hline

\end{tabular}%
}

\end{table}


This appendix provides evidence of each student's individual contribution to the group review.



% =========================================
% APPENDIX B - ARTICLE MATRIX
% =========================================

\newpage

\phantomsection
\addcontentsline{toc}{section}{Appendix B -- Article Matrix}

\begin{flushleft}
    {\fontsize{14pt}{17pt}\selectfont
    \textbf{APPENDIX B -- ARTICLE MATRIX}}
\end{flushleft}

\vspace{0.5cm}

The group should combine all individual article reviews into one matrix.


\vspace{0.5cm}








\begin{table}[H]

\caption{Group article matrix of reviewed research studies}
\label{tab:article-matrix}

\fontsize{10pt}{12pt}\selectfont

\resizebox{\textwidth}{!}{%
\begin{tabular}{|c|c|c|c|c|c|c|c|c|}
\hline

\textbf{Article} &
\textbf{Student} &
\textbf{Year} &
\textbf{Method} &
\textbf{Dataset} &
\textbf{Key Results} &
\textbf{Strength} &
\textbf{Limitation} &
\textbf{Research Gap}
\\
\hline

1 & Student 1 & & & & & & & 
\\
\hline

2 & Student 1 & & & & & & & 
\\
\hline

3 & Student 1 & & & & & & & 
\\
\hline

4 & Student 2 & & & & & & & 
\\
\hline

5 & Student 2 & & & & & & & 
\\
\hline

6 & Student 2 & & & & & & & 
\\
\hline

... & ... & & & & & & & 
\\
\hline

\end{tabular}%
}

\end{table}




\end{document}
